\labsection{4}{Кластеризация с помощью нейронной сети Кохонена}

\labpart{Цель работы}
Изучить архитектуру самоорганизующихся карт Кохонена, алгоритмы конкурентного обучения без учителя и
получить навыки кластеризации двумерных данных с визуальным анализом топологии сети.

\labpart{Теоретические сведения}
Раздел~\ref{sec:som} теоретической части: конкурентное обучение и правило WTA (\ref{eq:bmu}), правило
Кохонена (\ref{eq:som}), топологии \code{gridtop} и \code{hextop}, функция окрестности, фазы обучения
двухфазное правило обучения, показатели качества $QE$, $TE$ и U-матрица, реализация на Python
(подраздел~\ref{sec:py-som}).

\labpart{Постановка задачи}
Построить нейронную сеть Кохонена, разделяющую множество точек декартовой плоскости на заданное число
кластеров. Выборка генерируется программно: вокруг каждого центра из таблицы~\ref{tab:l4-var}
случайно располагаются 30--50 точек с разбросом (среднеквадратичным отклонением) $0,10$--$0,15$.

\begin{table}[!htb]
\caption{Варианты заданий к лабораторной работе №4}
\label{tab:l4-var}
\centering\small
\setlength{\tabcolsep}{4pt}
\begin{tabularx}{\textwidth}{c>{\raggedright}X|c>{\raggedright\arraybackslash}X}
\toprule
№ & Центры кластеров $(x; y)$ & № & Центры кластеров $(x; y)$ \\
\midrule
1 & $(-0,5; 0,5)$, $(0,5; 0,5)$, $(0; -0,6)$ & 6 & $(-0,6; 0,6)$, $(0,6; 0,6)$, $(0; 0)$, $(0; -0,7)$ \\
2 & $(-0,6; -0,4)$, $(0,6; -0,4)$, $(0; 0,7)$ & 7 & $(-0,5; -0,5)$, $(0,5; 0,5)$, $(-0,7; 0,6)$ \\
3 & $(\pm 0,5; 0,5)$, $(\pm 0,5; -0,5)$ & 8 & $(\pm 0,4; 0,4)$, $(0; -0,2)$, $(0,6; -0,6)$ \\
4 & $(\pm 0,7; 0)$, $(0; \pm 0,7)$ & 9 & $(-0,6; 0)$, $(0; 0,6)$, $(0,6; -0,6)$ \\
5 & $(-0,3; 0,3)$, $(0,4; -0,4)$, $(-0,6; -0,6)$ & 10 & $(\pm 0,5; 0)$, $(-0,5; 0,6)$, $(0,5; -0,6)$ \\
\bottomrule
\end{tabularx}
\end{table}

\labpart{Требования к программе}
Через интерфейс задаются топология решётки (\code{hextop} или \code{gridtop}), её размер (например,
1$\times$4, 2$\times$2, 3$\times$3, 5$\times$5) и число эпох обучения. На плоскости $[-1; 1]^2$
одновременно отображаются цветные маркеры точек выборки и текущие веса нейронов, соединённые линиями
решётки. Положение нейронов и связей перерисовывается каждые 10--20 эпох, так что видно, как
первоначально хаотичная сетка растягивается и покрывает облака точек. Скорость анимации регулируется
ползунком <<Задержка (мс)>>, процесс управляется кнопками <<Старт>>, <<Пауза>>, <<Сброс>> (в Jupyter~---
\code{ipywidgets.Dropdown}, \code{IntSlider}, \code{Button} и \code{FigureWidget} Plotly).

\labpart{Порядок выполнения работы}
\begin{steps}
\item Программно сформировать облака точек для кластеров варианта (\code{rng.normal(center, sigma, (n, 2))})
и отобразить их на плоскости.
\item Реализовать карту Кохонена на NumPy с числом нейронов, равным числу кластеров, и обучить её с
мультипликацией. Зафиксировать итоговые положения нейронов и сравнить их с центрами кластеров.
\item Провести серию экспериментов: сделать решётку избыточной (5$\times$5) и сравнить топологии
\code{hextop} и \code{gridtop}. Проанализировать, как веса распределяются внутри плотных облаков данных.
\item Для каждой конфигурации вычислить $QE$, $TE$ и число мёртвых нейронов (среднее по 5--10 запускам),
построить карту попаданий и U-матрицу для карты 5$\times$5.
\item Дополнительно: сравнить гауссову окрестность с двухфазным правилом, в котором радиус не сужается
ниже 1; сравнить результат с MiniSom и \code{KMeans}; оценить, сколько эпох нужно для упорядочивания
карты.
\end{steps}

\labpart{Пример выполнения}
Вариант~3: четыре кластера с центрами $(\pm 0,5; \pm 0,5)$, сгенерировано 47, 31, 48 и 46 точек (всего
172). Карта 2$\times$2 \code{hextop} с гауссовой окрестностью за 200 эпох поместила каждый нейрон в центр
своего облака с отклонением 0,005--0,012 от среднего кластера; $QE = 0,157$ совпадает с результатом
$k$-средних, разбиение совпадает с истинным (ARI~$= 1$). Результаты серии~--- в таблице~\ref{tab:l4-series}.

\begin{table}[!htb]
\caption{Показатели качества карт (200 эпох, среднее по 10 запускам)}
\label{tab:l4-series}
\centering\small
\begin{tabular}{lccccc}
\toprule
Конфигурация & Нейронов & $QE$ & $TE$ & Мёртвых нейронов & ARI \\
\midrule
1$\times$4 & 4 & 0,157 & 0,241 & 0 & 1,000 \\
2$\times$2 & 4 & 0,157 & 0,000 & 0 & 1,000 \\
3$\times$3 \code{gridtop} & 9 & 0,135 & 0,017 & 1,0 & 1,000 \\
3$\times$3 \code{hextop} & 9 & 0,134 & 0,000 & 0 & 1,000 \\
5$\times$5 \code{gridtop} & 25 & 0,077 & 0,108 & 3,7 & 1,000 \\
5$\times$5 \code{hextop} & 25 & 0,076 & 0,022 & 2,9 & 1,000 \\
\bottomrule
\end{tabular}
\end{table}

Избыточная карта 5$\times$5 (рисунок~\ref{fig:l4-train}) за 30--50 эпох упорядочивается в ровную
решётку, затем растягивается к облакам. Число нейронов на облако пропорционально числу его точек:
около 6 нейронов на облака из 46--48 точек и 3,4 на облако из 31 точки. Три--четыре нейрона остаются
между облаками без попаданий и образуют на U-матрице границы кластеров (рисунок~\ref{fig:l4-umatrix}).
Гексагональная решётка сохраняет топологию в 5 раз лучше прямоугольной ($TE = 0,022$ против $0,108$) при
одинаковой квантизационной ошибке. Двухфазное правило с несужающейся окрестностью (радиус 1) стягивает центроиды
карты 2$\times$2 к центру плоскости ($QE = 0,38$--$0,46$, рисунок~\ref{fig:som-rules}). Для упорядочивания
карте достаточно 20--50 эпох, после 100 эпох $QE$ почти не меняется.

\begin{figure}[!htb]
\centering
\includegraphics[width=\textwidth]{\labfig{4}{05_training_5x5_hextop.png}}
\caption{Мультипликация обучения избыточной карты 5$\times$5 \code{hextop}: эпохи 0, 50 и 200}
\label{fig:l4-train}
\end{figure}

\begin{figure}[!htb]
\centering
\includegraphics[width=0.5\textwidth]{\labfig{4}{06_hits_umatrix.png}}
\caption{Карта 5$\times$5: число попаданий точек на нейрон и U-матрица}
\label{fig:l4-umatrix}
\end{figure}

\labpart{Контрольные вопросы}
\begin{questions}
\item Каков биологический аналог структурной организации самоорганизующихся карт Кохонена?
\item Объясните математический принцип стратегии <<победитель получает всё>>. Какова связь с методом
$k$-средних?
\item Как функция окрестности $h_{ci}$ и топология решётки (\code{hextop}/\code{gridtop}) влияют на изменение
весов соседних нейронов?
\item Что происходит с картой при избыточном числе нейронов? Как в этом случае выделить кластеры?
\item Почему сеть Кохонена обучается без учителя? Откуда она извлекает информацию о структуре данных?
\item Что такое мёртвые нейроны? Что показывают квантизационная и топографическая ошибки, U-матрица?
\item Зачем радиус окрестности уменьшают в ходе обучения? Что будет, если он не сужается до нуля?
\end{questions}
